\section{Motivation}
\label{sec:intro:motivation}
% sections/1.Introduction/1.1.Motivation.tex

\begin{figure}[H]
\centering
\imageplaceholder{Photograph}{A wide shot of an HCMC public bus on a busy
street, capturing the network in context (passengers boarding, route
signage visible)}{Photograph an HCMC bus during peak hours (e.g.,
District~1, B\'{e}n Th\`{a}nh area). Save as
\texttt{images/screenshots/hcmc\_bus\_photo.jpg}.}
\caption{The HCMC public bus network --- subject of this work.}
\label{fig:intro:bus}
\end{figure}

Ho Chi Minh City (HCMC) operates a bus network of approximately 152 urban
routes, 90 of which are subsidized, served by a fleet of over 2,000 vehicles
managed by the Centre for Public Transport Management~\cite{ttgtcc2024}.
Each day, the network's electronic ticketing system records every boarding
event, capturing when and where passengers tap in, what fare class they
pay, and which trip and vehicle they board. Combined with GPS positioning
and GTFS schedule data, this creates a detailed operational record of the
entire network's activity.
 
Despite this data richness, operational decisions are still largely driven by
manual inspection of spreadsheet exports and static reports. Anomalies that
are clearly visible in the data: a stop that is chronically delayed in the
evening peak, a route with an excessive proportion of near-empty trips, or a
bus that consistently misses the connection window at a metro interchange, go
undetected until they escalate into service failures. There is no unified,
interactive tool that allows a route administrator, a company executive, or a
city planner to explore the network's spatial and temporal dimensions without
data science expertise.

\begin{figure}[H]
\centering
\imageplaceholder{Screenshot}{The MCPT public route directory web page at
\url{buyttgtcc.com.vn} (or the official MCPT route lookup), showing a
list/table of routes with no analytical or aggregation
capability}{Open the MCPT public route directory in a browser, navigate
to the main route-list view, capture a full-page screenshot at a
typical viewport (e.g., 1440~px wide). Save as
\texttt{images/screenshots/mcpt\_directory.png}.}
\caption{The MCPT public route directory~\cite{ttgtcc2024} ---
the only operational artefact for the network publicly accessible
at the time of writing. The directory is a navigational tool,
not an analytical one: routes can be looked up but no aggregation,
visualisation, or temporal analysis is offered.}
\label{fig:intro:mcpt}
\end{figure}
 
This gap is further sharpened by the opening of Metro Line~1 on 22 December 2024~\cite{metro_wiki}, which
introduced an intermodal dependency into the network for the first time.
The effectiveness of metro-based urban mobility now depends directly on the
punctuality of bus feeder services: a bus arriving five minutes late at an
interchange station translates directly into a missed metro departure for its
passengers. Monitoring this intermodal relationship at scale requires
integrating bus operational data with the metro timetable, a capability absent
from existing management tools.

From an operational standpoint, transport companies and city administrators
face concrete daily questions that this data could answer but currently does
not: Which routes are running near-empty trips that waste fuel and driver
hours? Are scheduled trips actually departing on time, and if not, at which
stops does delay accumulate? Which routes carry disproportionate numbers of
subsidized passengers, including students and elderly riders, and are fare subsidies
reaching those groups efficiently? How does ridership shift between peak and
off-peak hours, and is the current trip frequency matched to actual demand?
Without interactive visualization tools, answering any of these questions
requires manual extraction and cross-referencing of multiple raw tables,
a process that is slow, error-prone, and inaccessible to non-technical staff.

This thesis develops \textbf{BusLens}, an analytics and visualisation
platform that reconciles the three data streams behind a single
queryable surface and exposes the network's spatial, temporal, and
financial dimensions through six themed lenses ---
\emph{Demand \& Crowding}, \emph{School-Aware Demand},
\emph{Spatial \& Connectivity}, \emph{Financial \& Operations},
\emph{Network Health \& Performance}, and \emph{Routes}. Each lens
is one card in an interactive hub
(Chapter~\ref{chapter:prototype}); each card is built on the same
data pipeline (Chapter~\ref{chapter:proposed}) and obeys three
invariants that anchor the user experience: linked brushing,
story-subtitled charts, and a universal entity detail panel that
opens on every entity click.

Looking ahead, the 2025--2030 bus network restructuring plan targets a
complete reorganisation of the current routes into a tiered hierarchy of
179 routes (18 trunk, 73 branch, 88 feeder)~\cite{vnexpress2023}, alongside
a full transition to electronic ticketing by 2025 and a 20--25\% electric
fleet by 2030~\cite{vnexpress2023}. These reforms will substantially increase
the volume and complexity of operational data-making the absence of visual
analytics tools an increasingly critical problem.


